\documentclass[12pt]{article}
\usepackage[UTF8]{inputenc}
\usepackage{xeCJK}
\setCJKmainfont{Sarabun}
\usepackage{enumitem}
\usepackage{geometry}
\geometry{a4paper, margin=2.5cm}
\usepackage{setspace}
\onehalfspacing

\begin{document}

\begin{center}
{\Large\textbf{แสง: หักเห}}\\[6pt]
{\normalsize วิชาฟิสิกส์ ม.ปลาย บทที่ 22}\\[4pt]
{\footnotesize ทั้งหมด 25 ข้อ}
\end{center}
\hrule
\bigskip

\section*{คำถาม in class}
\begin{enumerate}[label=\arabic*., itemsep=0.5\baselineskip, topsep=0.5\baselineskip, leftmargin=*]
\item แสงเดินทางจากอากาศ (n = 1) เข้าสู่น้ำ (n = 4/3) ถ้ามุมตกเท่ากับ 30° มุมหักเหเท่ากับกี่องศา
\begin{enumerate}[label=\alph*)., itemsep=0.3\baselineskip, topsep=0.3\baselineskip]
\item ก\ 
\begin{minipage}[t]{0.9\linewidth}
\raggedright 22.0°
\end{minipage}
\item ข\ 
\begin{minipage}[t]{0.9\linewidth}
\raggedright 30°
\end{minipage}
\item ค\ 
\begin{minipage}[t]{0.9\linewidth}
\raggedright 44.6°
\end{minipage}
\item ง\ 
\begin{minipage}[t]{0.9\linewidth}
\raggedright 60°
\end{minipage}
\end{enumerate}
\vspace{-0.3\baselineskip}
\begin{small}
\textit{คำตอบ: ก}
\end{small}

\item เมื่อแสงเดินทางจากตัวกลางที่มีดัชนีหักเหน้อยไปยังตัวกลางที่มีดัชนีหักเหมาก แสงจะหักเหเข้าหาเส้นปกติหรือออกจากเส้นปกติ
\begin{enumerate}[label=\alph*)., itemsep=0.3\baselineskip, topsep=0.3\baselineskip]
\item ก\ 
\begin{minipage}[t]{0.9\linewidth}
\raggedright หักเหเข้าหาเส้นปกติ
\end{minipage}
\item ข\ 
\begin{minipage}[t]{0.9\linewidth}
\raggedright หักเหออกจากเส้นปกติ
\end{minipage}
\item ค\ 
\begin{minipage}[t]{0.9\linewidth}
\raggedright เดินทางตรงต่อไป
\end{minipage}
\item ง\ 
\begin{minipage}[t]{0.9\linewidth}
\raggedright สะท้อนกลับ
\end{minipage}
\end{enumerate}
\vspace{-0.3\baselineskip}
\begin{small}
\textit{คำตอบ: ก}
\end{small}

\item แสงเดินทางจากแก้ว (n = 1.5) ไปยังน้ำ (n = 4/3) ถ้ามุมตกเท่ากับ 45° มุมหักเหเท่ากับกี่องศา
\begin{enumerate}[label=\alph*)., itemsep=0.3\baselineskip, topsep=0.3\baselineskip]
\item ก\ 
\begin{minipage}[t]{0.9\linewidth}
\raggedright 33.8°
\end{minipage}
\item ข\ 
\begin{minipage}[t]{0.9\linewidth}
\raggedright 45°
\end{minipage}
\item ค\ 
\begin{minipage}[t]{0.9\linewidth}
\raggedright 56.3°
\end{minipage}
\item ง\ 
\begin{minipage}[t]{0.9\linewidth}
\raggedright ไม่สามารถหักเหได้
\end{minipage}
\end{enumerate}
\vspace{-0.3\baselineskip}
\begin{small}
\textit{คำตอบ: ค}
\end{small}

\item แสงเดินทางจากน้ำ (n = 4/3) ไปยังอากาศ (n = 1) มุมวิกฤต (critical angle) มีค่าเท่าใด
\begin{enumerate}[label=\alph*)., itemsep=0.3\baselineskip, topsep=0.3\baselineskip]
\item ก\ 
\begin{minipage}[t]{0.9\linewidth}
\raggedright 41.8°
\end{minipage}
\item ข\ 
\begin{minipage}[t]{0.9\linewidth}
\raggedright 48.8°
\end{minipage}
\item ค\ 
\begin{minipage}[t]{0.9\linewidth}
\raggedright 60°
\end{minipage}
\item ง\ 
\begin{minipage}[t]{0.9\linewidth}
\raggedright 74.6°
\end{minipage}
\end{enumerate}
\vspace{-0.3\baselineskip}
\begin{small}
\textit{คำตอบ: ข}
\end{small}

\item เมื่อแสงเดินทางจากแก้ว (n = 1.5) ไปยังอากาศ (n = 1) ที่มุมตก 50° จะเกิดปรากฏการณ์ใด
\begin{enumerate}[label=\alph*)., itemsep=0.3\baselineskip, topsep=0.3\baselineskip]
\item ก\ 
\begin{minipage}[t]{0.9\linewidth}
\raggedright หักเหปกติ
\end{minipage}
\item ข\ 
\begin{minipage}[t]{0.9\linewidth}
\raggedright หักเหบางส่วน
\end{minipage}
\item ค\ 
\begin{minipage}[t]{0.9\linewidth}
\raggedright สะท้อนกลับหมด
\end{minipage}
\item ง\ 
\begin{minipage}[t]{0.9\linewidth}
\raggedright สะท้อนบางส่วน
\end{minipage}
\end{enumerate}
\vspace{-0.3\baselineskip}
\begin{small}
\textit{คำตอบ: ค}
\end{small}

\item เลนส์เว้า (concave lens) ให้ภาพชนิดใด
\begin{enumerate}[label=\alph*)., itemsep=0.3\baselineskip, topsep=0.3\baselineskip]
\item ก\ 
\begin{minipage}[t]{0.9\linewidth}
\raggedright ภาพจริงหัวกลับ
\end{minipage}
\item ข\ 
\begin{minipage}[t]{0.9\linewidth}
\raggedright ภาพจริงหัวตั้ง
\end{minipage}
\item ค\ 
\begin{minipage}[t]{0.9\linewidth}
\raggedright ภาพเสมือนหัวตั้ง
\end{minipage}
\item ง\ 
\begin{minipage}[t]{0.9\linewidth}
\raggedright ภาพเสมือนหัวกลับ
\end{minipage}
\end{enumerate}
\vspace{-0.3\baselineskip}
\begin{small}
\textit{คำตอบ: ค}
\end{small}

\item เลนส์นูน (convex lens) มีความยาวโฟกัส 20 cm วางวัตถุห่างจากเลนส์ 40 cm ระยะภาพเท่ากับกี่ cm
\begin{enumerate}[label=\alph*)., itemsep=0.3\baselineskip, topsep=0.3\baselineskip]
\item ก\ 
\begin{minipage}[t]{0.9\linewidth}
\raggedright 20 cm
\end{minipage}
\item ข\ 
\begin{minipage}[t]{0.9\linewidth}
\raggedright 40 cm
\end{minipage}
\item ค\ 
\begin{minipage}[t]{0.9\linewidth}
\raggedright 60 cm
\end{minipage}
\item ง\ 
\begin{minipage}[t]{0.9\linewidth}
\raggedright 80 cm
\end{minipage}
\end{enumerate}
\vspace{-0.3\baselineskip}
\begin{small}
\textit{คำตอบ: ข}
\end{small}

\item ปรากฏการณ์ที่แสงไม่สามารถออกจากเพชรได้เมื่อตกกระทบผิวด้วยมุมมากเกินไป เรียกว่าอะไร
\begin{enumerate}[label=\alph*)., itemsep=0.3\baselineskip, topsep=0.3\baselineskip]
\item ก\ 
\begin{minipage}[t]{0.9\linewidth}
\raggedright การหักเห
\end{minipage}
\item ข\ 
\begin{minipage}[t]{0.9\linewidth}
\raggedright การสะท้อนกลับหมด
\end{minipage}
\item ค\ 
\begin{minipage}[t]{0.9\linewidth}
\raggedright การเลี้ยวเบน
\end{minipage}
\item ง\ 
\begin{minipage}[t]{0.9\linewidth}
\raggedright การแทรกสอด
\end{minipage}
\end{enumerate}
\vspace{-0.3\baselineskip}
\begin{small}
\textit{คำตอบ: ข}
\end{small}

\item แสงเดินทางจากอากาศเข้าน้ำด้วยมุมตก 60° ถ้าดัชนีหักเหของน้ำเท่ากับ 4/3 มุมหักเหเท่ากับกี่องศา
\begin{enumerate}[label=\alph*)., itemsep=0.3\baselineskip, topsep=0.3\baselineskip]
\item ก\ 
\begin{minipage}[t]{0.9\linewidth}
\raggedright 40.5°
\end{minipage}
\item ข\ 
\begin{minipage}[t]{0.9\linewidth}
\raggedright 45°
\end{minipage}
\item ค\ 
\begin{minipage}[t]{0.9\linewidth}
\raggedright 60°
\end{minipage}
\item ง\ 
\begin{minipage}[t]{0.9\linewidth}
\raggedright sin⁻¹(4/3) ไม่มีคำตอบ
\end{minipage}
\end{enumerate}
\vspace{-0.3\baselineskip}
\begin{small}
\textit{คำตอบ: ก}
\end{small}

\item เลนส์เว้ามีความยาวโฟกัส -15 cm วางวัตถุห่างจากเลนส์ 10 cm ระยะภาพเท่ากับกี่ cm
\begin{enumerate}[label=\alph*)., itemsep=0.3\baselineskip, topsep=0.3\baselineskip]
\item ก\ 
\begin{minipage}[t]{0.9\linewidth}
\raggedright -6 cm
\end{minipage}
\item ข\ 
\begin{minipage}[t]{0.9\linewidth}
\raggedright -15 cm
\end{minipage}
\item ค\ 
\begin{minipage}[t]{0.9\linewidth}
\raggedright -30 cm
\end{minipage}
\item ง\ 
\begin{minipage}[t]{0.9\linewidth}
\raggedright 6 cm
\end{minipage}
\end{enumerate}
\vspace{-0.3\baselineskip}
\begin{small}
\textit{คำตอบ: ก}
\end{small}

\item เลนส์นูน (convex lens) มีคุณสมบัติอย่างไร
\begin{enumerate}[label=\alph*)., itemsep=0.3\baselineskip, topsep=0.3\baselineskip]
\item ก\ 
\begin{minipage}[t]{0.9\linewidth}
\raggedright ขอบหนากลางบาง รวมแสง
\end{minipage}
\item ข\ 
\begin{minipage}[t]{0.9\linewidth}
\raggedright ขอบหนากลางบาง กระจายแสง
\end{minipage}
\item ค\ 
\begin{minipage}[t]{0.9\linewidth}
\raggedright ขอบบางกลางหนา รวมแสง
\end{minipage}
\item ง\ 
\begin{minipage}[t]{0.9\linewidth}
\raggedright ขอบบางกลางหนา กระจายแสง
\end{minipage}
\end{enumerate}
\vspace{-0.3\baselineskip}
\begin{small}
\textit{คำตอบ: ค}
\end{small}

\item มุมวิกฤตของแสงที่เดินทางจากแก้ว (n = 1.5) ไปยังอากาศมีค่าเท่าใด
\begin{enumerate}[label=\alph*)., itemsep=0.3\baselineskip, topsep=0.3\baselineskip]
\item ก\ 
\begin{minipage}[t]{0.9\linewidth}
\raggedright 30°
\end{minipage}
\item ข\ 
\begin{minipage}[t]{0.9\linewidth}
\raggedright 41.8°
\end{minipage}
\item ค\ 
\begin{minipage}[t]{0.9\linewidth}
\raggedright 48.8°
\end{minipage}
\item ง\ 
\begin{minipage}[t]{0.9\linewidth}
\raggedright 60°
\end{minipage}
\end{enumerate}
\vspace{-0.3\baselineskip}
\begin{small}
\textit{คำตอบ: ข}
\end{small}

\item เมื่อแสงตกกระทบผิวแบ่งระหว่างตัวกลางสองชนิดแล้วเกิดการสะท้อนกลับหมด ข้อใดจำเป็นต้องเป็นจริง
\begin{enumerate}[label=\alph*)., itemsep=0.3\baselineskip, topsep=0.3\baselineskip]
\item ก\ 
\begin{minipage}[t]{0.9\linewidth}
\raggedright แสงเดินทางจากตัวกลางที่หนักกว่าไปเบากว่า
\end{minipage}
\item ข\ 
\begin{minipage}[t]{0.9\linewidth}
\raggedright มุมตกเท่ากับมุมวิกฤต
\end{minipage}
\item ค\ 
\begin{minipage}[t]{0.9\linewidth}
\raggedright มุมตกมากกว่ามุมวิกฤต
\end{minipage}
\item ง\ 
\begin{minipage}[t]{0.9\linewidth}
\raggedright ตัวกลางทั้งสองมีดัชนีหักเหเท่ากัน
\end{minipage}
\end{enumerate}
\vspace{-0.3\baselineskip}
\begin{small}
\textit{คำตอบ: ค}
\end{small}

\item เลนส์นูนมีความยาวโฟกัส 10 cm วางวัตถุที่ระยะ 6 cm จากเลนส์ ภาพที่ได้เป็นภาพชนิดใด
\begin{enumerate}[label=\alph*)., itemsep=0.3\baselineskip, topsep=0.3\baselineskip]
\item ก\ 
\begin{minipage}[t]{0.9\linewidth}
\raggedright ภาพจริงหัวกลับ
\end{minipage}
\item ข\ 
\begin{minipage}[t]{0.9\linewidth}
\raggedright ภาพจริงหัวตั้ง
\end{minipage}
\item ค\ 
\begin{minipage}[t]{0.9\linewidth}
\raggedright ภาพเสมือนหัวตั้ง
\end{minipage}
\item ง\ 
\begin{minipage}[t]{0.9\linewidth}
\raggedright ภาพเสมือนหัวกลับ
\end{minipage}
\end{enumerate}
\vspace{-0.3\baselineskip}
\begin{small}
\textit{คำตอบ: ค}
\end{small}

\item กล้องจุลทรรศนอปติคัล (optical microscope) ใช้เลนส์ชนิดใดเป็นหลัก
\begin{enumerate}[label=\alph*)., itemsep=0.3\baselineskip, topsep=0.3\baselineskip]
\item ก\ 
\begin{minipage}[t]{0.9\linewidth}
\raggedright เลนส์เว้า
\end{minipage}
\item ข\ 
\begin{minipage}[t]{0.9\linewidth}
\raggedright เลนส์นูน
\end{minipage}
\item ค\ 
\begin{minipage}[t]{0.9\linewidth}
\raggedright เลนส์ราบ
\end{minipage}
\item ง\ 
\begin{minipage}[t]{0.9\linewidth}
\raggedright เลนส์ทรงกลม
\end{minipage}
\end{enumerate}
\vspace{-0.3\baselineskip}
\begin{small}
\textit{คำตอบ: ข}
\end{small}

\item แสงตกกระทบผิวน้ำ (n = 4/3) จากอากาศด้วยมุม 45° ถ้าวัดมุมหักเหได้ 32° ดัชนีหักเหของน้ำมีค่าเท่าใด
\begin{enumerate}[label=\alph*)., itemsep=0.3\baselineskip, topsep=0.3\baselineskip]
\item ก\ 
\begin{minipage}[t]{0.9\linewidth}
\raggedright 1.20
\end{minipage}
\item ข\ 
\begin{minipage}[t]{0.9\linewidth}
\raggedright 1.33
\end{minipage}
\item ค\ 
\begin{minipage}[t]{0.9\linewidth}
\raggedright 1.50
\end{minipage}
\item ง\ 
\begin{minipage}[t]{0.9\linewidth}
\raggedright 1.78
\end{minipage}
\end{enumerate}
\vspace{-0.3\baselineskip}
\begin{small}
\textit{คำตอบ: ข}
\end{small}

\item เส้นใยนำแสง (optical fiber) ทำงานบนหลักการใด
\begin{enumerate}[label=\alph*)., itemsep=0.3\baselineskip, topsep=0.3\baselineskip]
\item ก\ 
\begin{minipage}[t]{0.9\linewidth}
\raggedright การหักเหของแสง
\end{minipage}
\item ข\ 
\begin{minipage}[t]{0.9\linewidth}
\raggedright การสะท้อนกลับหมด
\end{minipage}
\item ค\ 
\begin{minipage}[t]{0.9\linewidth}
\raggedright การเลี้ยวเบน
\end{minipage}
\item ง\ 
\begin{minipage}[t]{0.9\linewidth}
\raggedright การสะท้อนปกติ
\end{minipage}
\end{enumerate}
\vspace{-0.3\baselineskip}
\begin{small}
\textit{คำตอบ: ข}
\end{small}

\item ทำไมตัวแท่งแก้วที่จุ่มลงในน้ำจึงดูเหมือนหักงอ
\begin{enumerate}[label=\alph*)., itemsep=0.3\baselineskip, topsep=0.3\baselineskip]
\item ก\ 
\begin{minipage}[t]{0.9\linewidth}
\raggedright การสะท้อนของแสง
\end{minipage}
\item ข\ 
\begin{minipage}[t]{0.9\linewidth}
\raggedright การหักเหของแสง
\end{minipage}
\item ค\ 
\begin{minipage}[t]{0.9\linewidth}
\raggedright การเลี้ยวเบนของแสง
\end{minipage}
\item ง\ 
\begin{minipage}[t]{0.9\linewidth}
\raggedright การกระจายของแสง
\end{minipage}
\end{enumerate}
\vspace{-0.3\baselineskip}
\begin{small}
\textit{คำตอบ: ข}
\end{small}

\item เลนส์นูนมีความยาวโฟกัส 12 cm ถ้าต้องการให้ได้ภาพจริงที่ขยาย 3 เท่า ต้องวางวัตถุที่ระยะกี่ cm จากเลนส์
\begin{enumerate}[label=\alph*)., itemsep=0.3\baselineskip, topsep=0.3\baselineskip]
\item ก\ 
\begin{minipage}[t]{0.9\linewidth}
\raggedright 8 cm
\end{minipage}
\item ข\ 
\begin{minipage}[t]{0.9\linewidth}
\raggedright 12 cm
\end{minipage}
\item ค\ 
\begin{minipage}[t]{0.9\linewidth}
\raggedright 16 cm
\end{minipage}
\item ง\ 
\begin{minipage}[t]{0.9\linewidth}
\raggedright 20 cm
\end{minipage}
\end{enumerate}
\vspace{-0.3\baselineskip}
\begin{small}
\textit{คำตอบ: ค}
\end{small}

\item เลนส์เว้ามีความยาวโฟกัส -20 cm วางวัตถุห่างจากเลนส์ 30 cm ภาพที่ได้มีการขยายเท่าใด
\begin{enumerate}[label=\alph*)., itemsep=0.3\baselineskip, topsep=0.3\baselineskip]
\item ก\ 
\begin{minipage}[t]{0.9\linewidth}
\raggedright 0.4 เท่า
\end{minipage}
\item ข\ 
\begin{minipage}[t]{0.9\linewidth}
\raggedright 0.67 เท่า
\end{minipage}
\item ค\ 
\begin{minipage}[t]{0.9\linewidth}
\raggedright 1.5 เท่า
\end{minipage}
\item ง\ 
\begin{minipage}[t]{0.9\linewidth}
\raggedright 2.5 เท่า
\end{minipage}
\end{enumerate}
\vspace{-0.3\baselineskip}
\begin{small}
\textit{คำตอบ: ก}
\end{small}

\item แสงเดินทางจากอากาศ (n = 1) เข้าสู่แก้ว (n = 1.5) ด้วยมุมตก 60° มุมหักเหเท่ากับกี่องศา
\begin{enumerate}[label=\alph*)., itemsep=0.3\baselineskip, topsep=0.3\baselineskip]
\item ก\ 
\begin{minipage}[t]{0.9\linewidth}
\raggedright 35.3°
\end{minipage}
\item ข\ 
\begin{minipage}[t]{0.9\linewidth}
\raggedright 40.0°
\end{minipage}
\item ค\ 
\begin{minipage}[t]{0.9\linewidth}
\raggedright 44.6°
\end{minipage}
\item ง\ 
\begin{minipage}[t]{0.9\linewidth}
\raggedright 60°
\end{minipage}
\end{enumerate}
\vspace{-0.3\baselineskip}
\begin{small}
\textit{คำตอบ: ก}
\end{small}

\item เลนส์นูนสองใบมี f₁ = 10 cm และ f₂ = 20 cm นำมาวางชิดกัน ความยาวโฟกัสรวมเป็นเท่าใด
\begin{enumerate}[label=\alph*)., itemsep=0.3\baselineskip, topsep=0.3\baselineskip]
\item ก\ 
\begin{minipage}[t]{0.9\linewidth}
\raggedright 6.67 cm
\end{minipage}
\item ข\ 
\begin{minipage}[t]{0.9\linewidth}
\raggedright 10 cm
\end{minipage}
\item ค\ 
\begin{minipage}[t]{0.9\linewidth}
\raggedright 15 cm
\end{minipage}
\item ง\ 
\begin{minipage}[t]{0.9\linewidth}
\raggedright 30 cm
\end{minipage}
\end{enumerate}
\vspace{-0.3\baselineskip}
\begin{small}
\textit{คำตอบ: ก}
\end{small}

\item เมื่อแสงเดินทางจากน้ำไปยังอากาศ ถ้ามุมตกเท่ากับ 30° จะเกิดการสะท้อนกลับหมดหรือไม่ (θc ≈ 48.8°)
\begin{enumerate}[label=\alph*)., itemsep=0.3\baselineskip, topsep=0.3\baselineskip]
\item ก\ 
\begin{minipage}[t]{0.9\linewidth}
\raggedright เกิด เพราะ 30° < 48.8°
\end{minipage}
\item ข\ 
\begin{minipage}[t]{0.9\linewidth}
\raggedright เกิด เพราะ 30° > 48.8°
\end{minipage}
\item ค\ 
\begin{minipage}[t]{0.9\linewidth}
\raggedright ไม่เกิด เพราะ 30° < 48.8°
\end{minipage}
\item ง\ 
\begin{minipage}[t]{0.9\linewidth}
\raggedright ไม่เกิด เพราะแสงเดินจากหนักไปเบา
\end{minipage}
\end{enumerate}
\vspace{-0.3\baselineskip}
\begin{small}
\textit{คำตอบ: ค}
\end{small}

\item เลนส์เว้า (concave lens) มีคุณสมบัติอย่างไร
\begin{enumerate}[label=\alph*)., itemsep=0.3\baselineskip, topsep=0.3\baselineskip]
\item ก\ 
\begin{minipage}[t]{0.9\linewidth}
\raggedright ขอบบางกลางหนา รวมแสง
\end{minipage}
\item ข\ 
\begin{minipage}[t]{0.9\linewidth}
\raggedright ขอบบางกลางหนา กระจายแสง
\end{minipage}
\item ค\ 
\begin{minipage}[t]{0.9\linewidth}
\raggedright ขอบหนากลางบาง รวมแสง
\end{minipage}
\item ง\ 
\begin{minipage}[t]{0.9\linewidth}
\raggedright ขอบหนากลางบาง กระจายแสง
\end{minipage}
\end{enumerate}
\vspace{-0.3\baselineskip}
\begin{small}
\textit{คำตอบ: ง}
\end{small}

\item เลนส์นูนมีความยาวโฟกัส 15 cm วางวัตถุที่ระยะ 20 cm จากเลนส์ ระยะภาพเท่ากับกี่ cm
\begin{enumerate}[label=\alph*)., itemsep=0.3\baselineskip, topsep=0.3\baselineskip]
\item ก\ 
\begin{minipage}[t]{0.9\linewidth}
\raggedright 30 cm
\end{minipage}
\item ข\ 
\begin{minipage}[t]{0.9\linewidth}
\raggedright 40 cm
\end{minipage}
\item ค\ 
\begin{minipage}[t]{0.9\linewidth}
\raggedright 60 cm
\end{minipage}
\item ง\ 
\begin{minipage}[t]{0.9\linewidth}
\raggedright 75 cm
\end{minipage}
\end{enumerate}
\vspace{-0.3\baselineskip}
\begin{small}
\textit{คำตอบ: ค}
\end{small}

\end{enumerate}
\end{document}